\chapter{Presupuesto}
\label{cap:presupuesto}

Este capítulo recoge el coste del proyecto. Se trata de una estimación: las horas se han asignado
por fases a partir del registro de trabajo del proyecto, y las tarifas corresponden a valores de
mercado para perfiles junior de ciencia de datos.

\section{Coste de personal}

El proyecto lo ha desarrollado una sola persona. La \cref{tab:presupuesto_personal} distribuye la
dedicación por fases.

\begin{table}[H]
  \centering
  \caption[Coste de personal]{Dedicación estimada por fases y coste de personal asociado.}
  \label{tab:presupuesto_personal}
  \begin{tabular}{lrrr}
    \toprule
    Fase & Horas & Tarifa (€/h) & Coste (€) \\
    \midrule
    Revisión bibliográfica y diseño experimental   &  45 & 30 & 1.350 \\
    Construcción del conjunto de datos             &  80 & 30 & 2.400 \\
    Métodos de referencia y condiciones A--C       &  70 & 30 & 2.100 \\
    Infraestructura federada y condiciones D--E    &  90 & 30 & 2.700 \\
    Ajuste de hiperparámetros y contraste estadístico & 55 & 30 & 1.650 \\
    Despliegue en la nube (MLOps)                  &  60 & 30 & 1.800 \\
    Análisis de negocio                            &  30 & 30 &   900 \\
    Redacción de la memoria                        &  70 & 30 & 2.100 \\
    \midrule
    \textbf{Total}                                 & \textbf{500} & & \textbf{15.000} \\
    \bottomrule
  \end{tabular}
\end{table}

\section{Coste de equipos}

Se ha empleado un ordenador portátil personal. Se imputa la amortización correspondiente al periodo
de uso del proyecto, seis meses, sobre una vida útil de cuatro años.

\begin{table}[H]
  \centering
  \caption[Coste de equipos]{Amortización de los equipos empleados.}
  \label{tab:presupuesto_equipos}
  \begin{tabular}{lrrr}
    \toprule
    Concepto & Precio (€) & Amortización & Coste imputado (€) \\
    \midrule
    Ordenador portátil & 1.400 & 6/48 meses & 175 \\
    \bottomrule
  \end{tabular}
\end{table}

\section{Coste de software y servicios en la nube}

La totalidad de las herramientas empleadas son de código abierto y sin coste de licencia: Python,
PyTorch, Flower, LightGBM, \emph{scikit-learn}, pandas, DVC, Hydra, Streamlit y la distribución
MiKTeX para la composición de esta memoria.

El único coste de servicios corresponde al despliegue distribuido en Microsoft Azure. El gasto real
del ejercicio fue inferior a 5\,€, gracias al uso de máquinas de la gama más económica y a la
disciplina de detenerlas al terminar cada sesión. Se recoge en la
\cref{tab:presupuesto_software} el coste real y, a título orientativo, el que tendría una
operación continuada del sistema durante un mes.

\begin{table}[H]
  \centering
  \caption[Coste de software y servicios]{Coste de software y de servicios en la nube.}
  \label{tab:presupuesto_software}
  \begin{tabular}{lr}
    \toprule
    Concepto & Coste (€) \\
    \midrule
    Licencias de software (código abierto) & 0 \\
    Microsoft Azure — experimentación realizada & 5 \\
    \midrule
    \textbf{Total imputado al proyecto} & \textbf{5} \\
    \addlinespace
    \emph{(Referencia: operación continua de 3 máquinas, 1 mes)} & \emph{$\approx$ 75} \\
    \bottomrule
  \end{tabular}
\end{table}

Merece la pena subrayar la asimetría entre este capítulo y el
\cref{cap:negocio}: el coste de cómputo del entrenamiento federado es marginal frente al beneficio
estimado. El obstáculo para adoptar un esquema como el estudiado no es económico ni tecnológico,
sino de gobierno del dato y de acuerdo entre las partes.

\section{Resumen del presupuesto}

\begin{table}[H]
  \centering
  \caption[Resumen del presupuesto]{Resumen del presupuesto del proyecto.}
  \label{tab:presupuesto_resumen}
  \begin{tabular}{lr}
    \toprule
    Partida & Coste (€) \\
    \midrule
    Personal                       & 15.000 \\
    Equipos (amortización)         &    175 \\
    Software y servicios en la nube &      5 \\
    \midrule
    \textbf{Subtotal}              & \textbf{15.180} \\
    Costes indirectos (10\,\%)     &  1.518 \\
    \midrule
    \textbf{Total}                 & \textbf{16.698} \\
    \bottomrule
  \end{tabular}
\end{table}
